% ===== PRÉAMBULE CANONIQUE — à coller VERBATIM en tête de CHAQUE .tex =====
\documentclass[11pt,a4paper]{article}
\usepackage[utf8]{inputenc}\usepackage[T1]{fontenc}\usepackage{lmodern}
\usepackage[french]{babel}
\usepackage{amsmath,amssymb,mathtools}\usepackage{icomma}
\usepackage{siunitx}\sisetup{locale=FR}  % \qty{}{}, \si{}, \ang{}
\usepackage{xcolor}\usepackage{colortbl}
\usepackage{tikz}\usepackage{pgfplots}\pgfplotsset{compat=1.18}
\usepackage[siunitx,europeanresistors]{circuitikz} % circuits (élec.)
\usepackage[most]{tcolorbox}\usepackage{enumitem}\usepackage{array,booktabs}
\usepackage[a4paper,margin=2.2cm]{geometry}
\usetikzlibrary{arrows.meta,calc,angles,quotes,positioning,patterns,%
  backgrounds,shapes.geometric,decorations.pathmorphing,decorations.markings,intersections}
\definecolor{primary}{RGB}{222,49,99}\definecolor{primarydark}{RGB}{150,22,55}
\definecolor{defcol}{RGB}{0,114,178}\definecolor{methcol}{RGB}{52,92,148}
\definecolor{excol}{RGB}{0,135,90}\definecolor{attncol}{RGB}{200,40,55}
\tcbset{boxbase/.style={enhanced,breakable,boxrule=0.9pt,arc=2.5pt,
  left=3mm,right=3mm,top=2.3mm,bottom=2.3mm,fonttitle=\bfseries,coltitle=white}}
% --- boîtes TUTORIEL (noms EXACTS, ne pas renommer) ---
\newtcolorbox{cle}[1][]{boxbase,colback=defcol!6,colframe=defcol,
  title={\textsc{À retenir}\ifx\relax#1\relax\relax\else\textmd{ : \itshape #1}\fi}}
\newtcolorbox{methode}[1][]{boxbase,colback=methcol!7,colframe=methcol,sharp corners,
  title={\textsc{Méthode}\ifx\relax#1\relax\relax\else\textmd{ --- \itshape #1}\fi}}
\newtcolorbox{exemplebox}[1][]{boxbase,colback=excol!5,colframe=excol,
  title={\textsc{Exemple}\ifx\relax#1\relax\relax\else\textmd{ : \itshape #1}\fi}}
\newtcolorbox{piege}[1][]{boxbase,colback=attncol!6,colframe=attncol,
  title={\textsc{Piège}\ifx\relax#1\relax\relax\else\textmd{ : \itshape #1}\fi}}
\newtcolorbox{remarque}[1][]{boxbase,colback=black!4,colframe=black!45,
  title={\textsc{Remarque}\ifx\relax#1\relax\relax\else\textmd{ : \itshape #1}\fi}}
% --- boîtes EXERCICE / CORRIGÉ (noms EXACTS, auto-comptées) ---
\newtcolorbox[auto counter]{exo}[1][]{boxbase,colback=primary!4,colframe=primary,
  title={\textsc{Exercice}~\thetcbcounter\ifx\relax#1\relax\relax\else\textmd{ --- \itshape #1}\fi}}
\newtcolorbox[auto counter]{corrige}[1][]{boxbase,colback=excol!4,colframe=excol,
  title={\textsc{Corrigé}~\thetcbcounter\ifx\relax#1\relax\relax\else\textmd{ --- \itshape #1}\fi}}
\newcommand{\vv}[1]{\vec{#1}}\newcommand{\ds}{\displaystyle}
\newcommand{\R}{\mathbb{R}}\newcommand{\N}{\mathbb{N}}\newcommand{\Z}{\mathbb{Z}}
% (un \usepackage{fancyhdr}+en-tête est possible ; il est ignoré côté web)
% =========================================================================
\begin{document}
\begin{exo}[deux champs perpendiculaires de même norme]
En un point $M$, deux sources créent deux champs \textbf{perpendiculaires}, l'un horizontal
$\vv{B}_1$, l'autre vertical $\vv{B}_2$, de \textbf{même intensité} $B_1=B_2=\qty{1.0e-4}{\tesla}$.
\begin{center}
\begin{tikzpicture}[>=Stealth,scale=1.0]
  \fill[black] (0,0) circle (1.3pt) node[below left,font=\footnotesize]{$M$};
  \draw[blue!70!black,line width=1.4pt,->] (0,0) -- (2,0) node[right]{$\vv{B}_1$};
  \draw[blue!70!black,line width=1.4pt,->] (0,0) -- (0,2) node[above]{$\vv{B}_2$};
\end{tikzpicture}
\end{center}
\begin{enumerate}[label=\alph*)]
  \item Calcule l'intensité du champ résultant $\vv{B}_{\text{tot}}$.
  \item Quel angle $\theta$ fait-il avec $\vv{B}_1$ ?
\end{enumerate}
\end{exo}

\begin{exo}[superposition orthogonale : Pythagore]
Deux champs perpendiculaires se superposent en $M$ : $B_1=\qty{3e-4}{\tesla}$ (horizontal) et
$B_2=\qty{4e-4}{\tesla}$ (vertical).
\begin{enumerate}[label=\alph*)]
  \item Détermine l'intensité du champ résultant.
  \item Détermine l'angle $\theta$ que fait le champ résultant avec le champ horizontal $\vv{B}_1$
        (on donne $\tan^{-1}(4/3)\approx\ang{53}$).
\end{enumerate}
\end{exo}

\begin{exo}[fil horizontal : champ au-dessus et au-dessous]
Un fil rectiligne \textbf{horizontal}, dans le plan de la feuille, est parcouru par un courant $I$
orienté \textbf{vers la droite}. On considère un point $P$ situé \emph{au-dessus} du fil et un point
$Q$ situé \emph{au-dessous}.
\begin{center}
\begin{tikzpicture}[>=Stealth,scale=1.0]
  \draw[red!80!black,line width=2pt] (-2.4,0) -- (2.4,0);
  \draw[red!80!black,->,line width=2pt] (0.6,0) -- (1.4,0);
  \node[red!80!black,above] at (1.9,0){$I$};
  \fill[black] (-0.6,0.9) circle (1.4pt) node[above,font=\footnotesize]{$P$};
  \fill[black] (-0.6,-0.9) circle (1.4pt) node[below,font=\footnotesize]{$Q$};
\end{tikzpicture}
\end{center}
\begin{enumerate}[label=\alph*)]
  \item Au point $P$ (au-dessus), $\vv{B}$ sort-il ($\odot$) ou entre-t-il ($\otimes$) dans la feuille ?
  \item Même question au point $Q$ (au-dessous).
\end{enumerate}
\end{exo}

\begin{exo}[deux fils identiques : champ au milieu]
Deux fils parallèles, perpendiculaires à la feuille, portent le \textbf{même courant} de \textbf{même
sens} (tous deux \textbf{sortants}, $\odot$). Le point $M$ est situé exactement au \textbf{milieu} des
deux fils ; chacun y crée un champ d'intensité $\qty{4e-5}{\tesla}$.
\begin{center}
\begin{tikzpicture}[>=Stealth,scale=1.0]
  \foreach \x/\n in {-2/1,2/2}{
     \draw[red!80!black,line width=1.4pt] (\x,0) circle (0.16);
     \fill[red!80!black] (\x,0) circle (0.05);
     \node[red!80!black,below,font=\footnotesize] at (\x,-0.3){fil \n};
  }
  \fill[black] (0,0) circle (1.5pt) node[below,font=\footnotesize]{$M$};
\end{tikzpicture}
\end{center}
\begin{enumerate}[label=\alph*)]
  \item En appliquant la règle de la main droite à chaque fil, indique la direction et le sens des
        champs $\vv{B}_1$ et $\vv{B}_2$ créés en $M$.
  \item En déduis l'intensité du champ résultant en $M$.
\end{enumerate}
\end{exo}

\end{document}
